% IntegriRef: Multi-Layer Cross-Domain Reference Integrity Verification
% LaTeX manuscript — compile with pdflatex or xelatex
\documentclass[sigconf,nonacm]{acmart}

\usepackage{booktabs}
\usepackage{multirow}
\usepackage{amsmath}
\usepackage{algorithm}
\usepackage{algorithmic}
\usepackage{graphicx}
\usepackage{subcaption}
\usepackage{xcolor}
\usepackage{listings}
\usepackage{hyperref}
\usepackage{enumitem}

% Listings style
\lstset{
  basicstyle=\ttfamily\small,
  breaklines=true,
  frame=single,
  language=Python,
  keywordstyle=\color{blue},
  commentstyle=\color{gray},
}

\begin{document}

\title{IntegriRef: A Five-Layer Bayesian Framework for\\Cross-Domain Reference Integrity Verification}

\author{Anonymous}
\affiliation{%
  \institution{Anonymous Institution}
  \country{Anonymous}
}
\email{anonymous@institution.edu}

\begin{abstract}
The proliferation of AI-generated and paper-mill-produced academic content has created an urgent need for automated reference integrity verification.
Existing tools address isolated aspects of this problem---citation existence checking, semantic alignment, or graph anomaly detection---but none integrates these signals into a unified, calibrated risk assessment.
We present \textsc{IntegriRef}, a five-layer verification pipeline that combines registry-based existence checking (L0), citation intent classification (L1), natural language inference for claim--source alignment (L2), citation graph anomaly detection (L3), and Bayesian risk scoring with domain-specific priors (L4).
The system queries 62 registry adapters across six domains (academic, patents, legal, government, financial, standards) through an intelligent discovery layer that auto-routes queries based on identifier type.

Our L2 semantic NLI module achieves 93.5\% alignment accuracy on SciFact using fine-tuned DeBERTa-v3, exceeding the 89.1\% inter-annotator agreement ceiling.
The L1 intent classifier reaches 86.0\% macro F1 on SciCite, within 3 points of published SOTA despite using a 110M-parameter model.
On the IntegriRef-Bench benchmark of 58 golden test cases spanning 18 integrity signals, the full L0--L4 pipeline achieves 73.9\% retraction recall (17/23), 100\% chimera detection, 0\% false positive rate on 18 legitimate papers, and 100\% hallucination recall (14/14 fabricated references detected) after optimization of the phantom DOI detection and Bayesian kill-shot override mechanisms.
In a head-to-head comparison, IntegriRef outperforms CheckIfExist~\cite{alkhalaf2025} on all recall metrics while maintaining 0\% FPR, vs.\ CheckIfExist's 11.1\% FPR and inability to produce high-confidence detections.
A probabilistic early-stopping cascade achieves a 67.2\% L1--L3 skip rate with zero accuracy regressions.
Throughput optimizations---response caching, concurrent batch processing, ONNX INT8 quantization, and early-stop parallel queries---yield a 2.6$\times$ speedup for cold queries and 33,333 references/second for cached lookups.

We release IntegriRef as an open-source Python library with a REST API, together with IntegriRef-Bench, a novel multi-signal benchmark dataset covering retracted papers, LLM hallucinations, metadata chimeras, citation rings, tortured phrases, and statistical anomalies, curated from documented fraud cases in the academic integrity literature.
\end{abstract}

\keywords{reference integrity, citation verification, Bayesian risk scoring, citation graph analysis, NLI, paper mills, retraction detection}

\maketitle

% ========================================================================
\section{Introduction}
\label{sec:intro}

The integrity of scholarly references is fundamental to the trustworthiness of scientific communication.
A citation serves as a verifiable link between a claim and its evidentiary basis; when this link is broken---through fabricated references, phantom DOIs, metadata mismatches, or semantic misrepresentation---the entire chain of scientific reasoning is compromised.

Recent estimates suggest that problematic papers constitute approximately 2--10\% of published literature depending on the domain, with cancer research exhibiting rates as high as 9.87\% \cite{scancar2025}.
The Retraction Watch database has catalogued over 40,000 retracted papers, with the annual retraction rate accelerating from approximately 400/year in 2010 to over 10,000 in 2023 \cite{vannoorden2023}.
The IOP mass retraction of 350+ papers in 2024 for citation manipulation \cite{else2024iop}, and the Hindawi retraction of 511 special-issue papers in 2022, demonstrate that citation fraud operates at industrial scale.

The emergence of large language models (LLMs) has further exacerbated this problem by enabling the generation of plausible but entirely fabricated references.
Unlike traditional plagiarism, where copied text can be detected through string matching, LLM-generated references feature syntactically correct DOI formats, plausible author names, and realistic publication venues---yet point to nonexistent works.
Studies have shown that ChatGPT fabricates references in 30--70\% of generated academic text \cite{alkaissi2023}, with fabricated DOIs that follow valid publisher prefix patterns (e.g., \texttt{10.1038/} for Nature).

Paper mills represent an even more insidious threat.
Investigations have identified networks producing thousands of fraudulent papers: the ``Stock Photo Mill'' (121 papers with shared stock images), the ``Tadpole Paper Mill'' (400+ papers), and the Russian 123mi.ru mill (4,000+ papers) \cite{byrne2022}.
These mills employ sophisticated evasion techniques including tortured phrases---circumlocution paraphrases generated by synonym substitution to evade plagiarism detection (e.g., ``profound learning'' for ``deep learning'', ``property extrication'' for ``feature extraction'') \cite{cabanac2021}---and fabricated citation networks designed to artificially inflate impact factors.

Existing verification tools address fragments of this challenge:

\begin{itemize}[leftmargin=*]
    \item \textbf{Registry services} (Crossref, OpenAlex, PubMed) verify identifier existence but cannot assess whether a citation accurately represents the cited work's findings, nor do they detect subtle metadata fabrication.
    \item \textbf{Citation context tools} (scite.ai \cite{scite2019}) provide citation context classification but are limited to a single registry domain. Independent evaluation by Bakker et al.\ (2023) found scite.ai achieves only 29\% macro F1 due to severe mentioning-class bias.
    \item \textbf{Paper mill detectors} (Papermill Alarm, tortured-phrase detectors \cite{cabanac2021}) target specific fabrication signals but operate in isolation from broader verification pipelines.
    \item \textbf{Fact verification} (CiteTrue, RefChecker) combines existence checking with NLI but is limited to single registry domains and lacks calibrated risk quantification.
    \item \textbf{Statistical auditing} (GRIM test \cite{brown2017grim}, statcheck \cite{nuijten2016}) detects impossible reported statistics but has never been integrated into reference-level verification.
\end{itemize}

We identify three key limitations of the current landscape:

\begin{enumerate}
    \item \textbf{Signal fragmentation}: Verification signals (existence, intent, semantics, graph structure, text quality, statistics) are assessed independently, with no principled method for combining them into a unified risk score.
    \item \textbf{Domain silos}: Tools are built for specific identifier types (DOIs, PMIDs) and registries (Crossref, PubMed), leaving patents, legal citations, government reports, and financial filings unverified.
    \item \textbf{Uncalibrated confidence}: Binary ``found/not found'' verdicts provide no quantification of uncertainty, making it impossible to prioritize review effort across a large reference list.
\end{enumerate}

We present \textsc{IntegriRef}, a system that addresses all three limitations through a five-layer pipeline architecture integrating 18 integrity signals into a Bayesian risk model.
Our contributions are:

\begin{enumerate}
    \item \textbf{A unified five-layer architecture} that integrates existence, intent, semantic, graph, text quality, and statistical signals into a single verification pipeline with conditional layer escalation, querying 62 registry adapters across 6 domains.

    \item \textbf{A Bayesian risk scoring framework} with 18 literature-derived likelihood ratios and 8 domain-specific priors, producing calibrated posterior probabilities that enable principled prioritization of editorial review.

    \item \textbf{State-of-the-art semantic verification}: L2 NLI achieves 93.5\% accuracy on SciFact (exceeding human IAA of 89.1\%), and L1 intent classification reaches 86.0\% macro F1 on SciCite within 3 points of SOTA.

    \item \textbf{IntegriRef-Bench}, a novel benchmark dataset of 464 curated test cases covering all 18 integrity signals, drawn from documented fraud cases including 28 famous retractions, citation cartels, tortured-phrase corpora, and known statistical fabrication.

    \item \textbf{Production-ready implementation} with dual-layer caching, circuit breakers, ONNX INT8 quantization, async I/O, and a REST API, achieving 33,333 refs/s for cached queries.
\end{enumerate}

% ========================================================================
\section{Related Work}
\label{sec:related}

\subsection{Reference Existence Verification}

The foundational layer of reference verification checks whether a cited work actually exists.
DOI resolution via Crossref \cite{crossref2024} and DataCite provides high-confidence existence checks for documents with registered identifiers, covering over 150 million works.
OpenAlex \cite{priem2022openalex} aggregates metadata from multiple sources into a unified index of 250 million+ works, providing broader coverage but requiring fuzzy matching for title-based searches.
PubMed indexes 36 million biomedical citations, Semantic Scholar \cite{kinney2023semantic} covers 200 million papers with citation context, and DBLP provides curated computer science bibliographies.
However, these services operate independently: no existing tool orchestrates queries across multiple registries with intelligent routing based on identifier type, nor do they cover non-academic domains (patents, legal, government, financial, standards).

\subsection{Citation Context Analysis}

Cohan et al.\ \cite{cohan2019structural} introduced structural scaffolds for citation intent classification on the SciCite dataset, achieving 86.7\% F1 with SciBERT.
Jurgens et al.\ \cite{jurgens2018measuring} proposed citation frames for tracking how citations function within scientific discourse.
ImpactCite \cite{impactcite2020} achieved the current SOTA of 88.93\% macro F1 using XLNet-Large with multi-task learning.
Scite.ai \cite{scite2019} productized citation context classification but operates as a closed platform; independent evaluation found only 29\% macro F1 due to systematic bias toward the majority mentioning class \cite{bakker2023}.
The ACL-ARC dataset \cite{jurgens2018measuring} provides a finer-grained taxonomy (Background, Motivation, Uses, Extends, Comparison, Future) but with only 1,941 annotated citations.

\subsection{Semantic Verification via NLI}

Natural language inference has been applied to fact verification through systems like FEVER \cite{thorne2018fever} (185,000 claims against Wikipedia) and SciFact \cite{wadden2020fact} (1,409 scientific claims with expert-annotated evidence).
Wright et al.\ \cite{wright2022generating} demonstrated that cross-encoder NLI models can achieve high accuracy on scientific claim verification.
MultiVerS \cite{multivers2022} achieved 72.5\% F1 on end-to-end SciFact (retrieval + label prediction), while label-only prediction with fine-tuned DeBERTa reaches approximately 88\% F1 \cite{kosprdic2024}.
However, these systems have not been integrated into reference verification pipelines, and no prior work has established thresholds specifically for detecting citation-level misrepresentation (as opposed to standalone fact-checking).

\subsection{Citation Graph Anomaly Detection}

Fister et al.\ \cite{fister2016toward} proposed network-based approaches for detecting citation cartels, identifying mutual citation clusters through community detection on directed citation networks.
Kojaku et al.\ \cite{kojaku2021} developed the CIDRE algorithm for detecting mutual citation rings, demonstrating effectiveness on known cartel cases including the Brazilian GMR cartel (4 journals, 2013) and the IJNSNS/Ji-Huan He case.
Campanario \& Coslado \cite{campanario2011benford} applied Benford's law to citation count distributions, finding that deviations from the expected first-digit distribution correlate with citation manipulation.
Van Noorden \cite{vannoorden2013} documented the Brazilian citation stacking scandal that led Thompson Reuters to ban four journals.
These approaches have been studied independently but never integrated into a multi-signal verification pipeline with other evidence layers.
Xia et al.~\cite{xia2023action} propose ACTION, a framework using non-negative matrix factorization and network representation learning on heterogeneous academic networks to detect irrelevant citations.
Ma et al.~\cite{ma2025graphad} survey deep graph anomaly detection methods including GNN-based approaches applicable to citation network analysis.

\subsection{LLM Citation Hallucination}

The emergence of LLMs as writing tools has created a new category of reference integrity threats.
Alkaissi \& McFarlane~\cite{alkaissi2023} found that ChatGPT fabricates references in 30--70\% of generated academic text.
More recent work by Gravel et al.~\cite{gravel2024} shows that even GPT-4o fabricates 19.9\% of citations, with 45.4\% of ``real'' citations containing errors---DOIs exhibiting the highest error rate.
These fabricated references typically feature syntactically valid DOI formats, plausible author names, and realistic venues, making detection non-trivial.

Alkhalaf et al.~\cite{alkhalaf2025} present CheckIfExist, an open-source tool implementing cascading validation against CrossRef, Semantic Scholar, and OpenAlex---the most directly comparable system to IntegriRef's L0 layer, though it lacks semantic verification, graph analysis, and Bayesian risk scoring.
Heydari et al.~\cite{heydari2025} analyze citation planting through AI-generated papers on preprint servers and purchased citations from boosting services, documenting the scale of citation manipulation infrastructure.

Wu et al.~\cite{wu2024llmdetect} survey LLM-generated text detection methods broadly, covering watermarking, statistical, and neural approaches.
However, reference-level hallucination detection (as opposed to general text detection) requires domain-specific signals---phantom DOI patterns, author-title chimeras, temporal impossibilities---that general detectors cannot capture.

\subsection{Paper Mill Detection and Statistical Auditing}

Cabanac et al.\ \cite{cabanac2021} pioneered tortured-phrase detection, identifying 384+ characteristic circumlocution patterns produced by paper mills.
Their work revealed that journals such as \emph{Microprocessors and Microsystems} contained approximately 500 papers with tortured phrases, leading to mass retractions.
The GRIM test \cite{brown2017grim} checks whether reported means are mathematically possible given the sample size---for example, a mean of 3.53 with $N=25$ is impossible since $25 \times 3.53 = 88.25$ is not an integer.
Application of GRIM to Brian Wansink's ``pizza papers'' contributed to 18 retractions \cite{vanderzee2017}.
Statcheck \cite{nuijten2016} recomputes p-values from APA-formatted test statistics, finding that approximately 50\% of psychology papers contain at least one inconsistency.
Wilhite \& Fong \cite{wilhite2012} documented coercive citation practices and sneaked references---citations added to metadata but absent from the paper body.
None of these signals has been combined within a Bayesian framework for calibrated risk assessment.

% ========================================================================
\section{System Architecture}
\label{sec:architecture}

\subsection{Overview}

IntegriRef follows a layered pipeline architecture where each layer produces structured signals that feed into the final Bayesian risk scorer (Figure~\ref{fig:architecture}).
Layers can be selectively enabled or disabled, and the pipeline supports \emph{conditional escalation}---L2 (NLI verification) is only invoked when L0 flags ambiguity (composite confidence in the range 0.30--0.70), avoiding expensive model inference for clear-cut cases.

\begin{figure}[t]
\centering
\fbox{\parbox{0.45\textwidth}{\small
\textbf{Input}: Reference metadata + citing context \\[3pt]
$\downarrow$ \\
\textbf{L0}: Registry Verification (62 adapters, 6 domains) \\
\quad $\rightarrow$ existence, metadata match, retraction, hallucination \\[2pt]
$\downarrow$ \\
\textbf{L1}: Citation Intent Classification (SciBERT / heuristic) \\
\quad $\rightarrow$ intent label + confidence (5 classes) \\[2pt]
$\downarrow$ (conditional on L0 ambiguity) \\
\textbf{L2}: Semantic NLI Verification (DeBERTa-v3) \\
\quad $\rightarrow$ entailment / contradiction / neutral scores \\[2pt]
$\downarrow$ \\
\textbf{L3}: Citation Graph Analysis (OpenAlex 2-hop) \\
\quad $\rightarrow$ 7 anomaly detectors + graph health score \\[2pt]
$\downarrow$ \\
\textbf{L0$^\star$}: Text \& Statistical Signals \\
\quad $\rightarrow$ tortured phrases, email risk, GRIM, statcheck \\[2pt]
$\downarrow$ \\
\textbf{L4}: Bayesian Risk Scoring (18 signals, 8 priors) \\
\quad $\rightarrow$ posterior probability + risk tier (4 levels) \\[3pt]
\textbf{Output}: \texttt{PipelineReport} with provenance trail
}}
\caption{IntegriRef five-layer pipeline architecture. L0$^\star$ denotes text/statistical signals computed alongside L0 when full text is available. Conditional edges are triggered by upstream ambiguity.}
\label{fig:architecture}
\end{figure}

\subsection{Data Model: ICEntity}

The core data model centers on \texttt{ICEntity} (Integrity-Checked Entity), a unified representation that normalizes metadata from all 62 registry sources.
ICEntity supports 12 entity types: \texttt{paper}, \texttt{patent}, \texttt{case} (legal), \texttt{statute}, \texttt{standard}, \texttt{report}, \texttt{filing} (financial), \texttt{dataset}, \texttt{book}, \texttt{thesis}, \texttt{preprint}, and \texttt{other}.
Each entity carries:

\begin{itemize}[leftmargin=*]
    \item \textbf{Identity}: title, authors, year, identifiers (DOI, PMID, arXiv ID, patent number, case citation, CELEX, RFC, CVE, ISBN, CIK, ISO number)
    \item \textbf{Bibliographic}: venue, volume, pages, publisher, abstract (up to 500 characters)
    \item \textbf{Provenance}: source registry, query timestamp, raw response hash
    \item \textbf{Integrity metadata}: retraction status (\texttt{is\_retracted}), update history, expression of concern flags
\end{itemize}

\subsection{L0: Registry Verification Engine}
\label{sec:l0}

The L0 layer implements a multi-phase verification workflow that cascades from fast identifier lookup to comprehensive search:

\subsubsection{Phase 1: Identifier Detection and Parallel Lookup}

The \texttt{RegistryDiscovery} module auto-detects 15 identifier types using compiled regex patterns:

\begin{itemize}[leftmargin=*]
    \item \textbf{Academic}: DOI (\texttt{10.\textbackslash d\{4,9\}/}), arXiv (\texttt{\textbackslash d\{4\}.\textbackslash d\{4,5\}}), PMID, ORCID
    \item \textbf{Legal}: CELEX (EU law), case citations (e.g., \texttt{570 U.S. 529})
    \item \textbf{Patents}: USPTO (\texttt{US\textbackslash d\{7,8\}}), EPO (\texttt{EP\textbackslash d+}), WIPO (\texttt{WO\textbackslash d+})
    \item \textbf{Standards}: RFC (\texttt{RFC \textbackslash d+}), ISO (\texttt{ISO \textbackslash d+}), CVE
    \item \textbf{Financial}: CIK, EDGAR accession numbers
    \item \textbf{Books}: ISBN-10/13
\end{itemize}

Once identified, parallel registry queries are dispatched via \texttt{ThreadPoolExecutor} (default: 8 workers) or \texttt{asyncio.gather} for the async variant.
An \textbf{early-stop} mechanism cancels remaining futures after 3 confirmed matches, reducing average query time for well-indexed references while ensuring retraction metadata from slower registries (e.g., OpenAlex) is not missed.

\subsubsection{Phase 2: Search Fallback}

If no identifier matches, the system falls back to title/author/year search across domain-appropriate registries.
Registry selection is guided by heuristic domain classification based on venue name, identifier prefix, and author affiliation patterns.

\subsubsection{Phase 3: Composite Confidence Scoring}

For each registry response, a composite match score is computed:

\begin{equation}
C = 0.35 \cdot e + 0.25 \cdot s_t + 0.15 \cdot s_a + 0.10 \cdot s_v + 0.15 \cdot s_i
\label{eq:composite}
\end{equation}

\noindent where $e \in \{0, 1\}$ is existence, $s_t$ is title similarity (RapidFuzz token set ratio), $s_a$ is author similarity (normalized Jaccard over surname tokens), $s_v$ is venue similarity, and $s_i$ is context intent score.
The best-matching source determines the overall verdict:

\begin{itemize}[leftmargin=*]
    \item $C \geq 0.85$: \texttt{PASS} (high-confidence match)
    \item $0.60 \leq C < 0.85$: \texttt{WARN} (ambiguous, triggers L2 escalation)
    \item $C < 0.60$: \texttt{FAIL} (likely fabricated or severely mismatched)
\end{itemize}

\subsubsection{Phase 4: Retraction Detection}

Retraction status is checked through two independent mechanisms:
\begin{enumerate}
    \item \textbf{Crossref}: The \texttt{update-to} field in Crossref metadata links retraction notices to original papers. We follow these links to determine retraction status and extract retraction dates and reasons.
    \item \textbf{OpenAlex}: The \texttt{is\_retracted} boolean field provides a direct retraction flag aggregated from multiple sources including Retraction Watch.
\end{enumerate}

The dual-source approach is critical because Crossref stores retraction metadata on the \emph{retraction notice} (not the original paper), while OpenAlex stores it on the original paper.
Our early-stop threshold of 3 (rather than 2) ensures OpenAlex results are collected even when faster registries (Crossref, Semantic Scholar) respond first.

\subsubsection{Phase 5: Hallucination Detection}

The \texttt{HallucinationDetector} module applies 8 pattern-based checks for AI-generated reference artifacts:

\begin{enumerate}[leftmargin=*]
    \item \textbf{Phantom DOI}: DOI has valid format (\texttt{10.\textbackslash d\{4,\}/...}) but does not resolve via Crossref ID lookup (LR = 25.0). Critically, this signal fires \emph{regardless of whether title-based search finds similar papers}---a valid-format DOI that does not resolve is ground-truth evidence of fabrication, even if a similarly-titled real paper exists.
    \item \textbf{Phantom arXiv}: arXiv ID has valid format but does not exist. Same logic as phantom DOI.
    \item \textbf{No registry match}: Reference not found in any registry by ID. When search finds similar (but unconfirmed) papers, the signal fires at reduced confidence.
    \item \textbf{Author-title chimera}: API finds the author's real papers but with entirely different titles---a hallmark of LLM hallucination where a real researcher name is combined with an invented title.
    \item \textbf{Future date}: Publication year exceeds the current year.
    \item \textbf{Temporal impossibility}: Venue did not exist at the claimed year (e.g., NeurIPS in 1990).
    \item \textbf{Format anomaly}: Suspiciously uniform formatting across a reference list (e.g., $>$70\% use ``Title: Subtitle'' pattern), detected via batch analysis.
\end{enumerate}

\noindent\textbf{Kill-shot override.}
When the hallucination score exceeds a threshold ($\geq 50/100$), the pipeline forces the Bayesian risk tier to at least HIGH, preventing downstream Bayesian softening from masking strong phantom ID evidence.
Additionally, when a DOI or arXiv ID fails to resolve, the search fallback (Phase~2) enforces strict first-author surname matching, preventing false ``found'' results from similarly-titled papers.

The detector achieves 336,136 refs/s throughput, adding negligible latency to the pipeline.

\subsubsection{Registry Adapter Coverage}

Table~\ref{tab:registries} summarizes the 62 adapters across 6 domains.
Each adapter inherits from \texttt{RegistryAdapter}, which provides rate limiting, exponential backoff retries, and circuit breaker protection (3-state: closed/open/half-open with sliding window failure counting).

\begin{table}[t]
\caption{Registry adapter coverage by domain. Each adapter implements \texttt{query\_by\_id()}, \texttt{search()}, and \texttt{verify\_exists()} against a specific data source.}
\label{tab:registries}
\centering
\small
\begin{tabular}{lrl}
\toprule
\textbf{Domain} & \textbf{Count} & \textbf{Key Adapters} \\
\midrule
Academic & 35 & Crossref, S2, OpenAlex, PubMed, arXiv, \\
         &    & DBLP, CORE, Europe PMC, bioRxiv, \\
         &    & Zenodo, NASA ADS, INSPIRE-HEP, \\
         &    & CiNii, J-Stage, HAL, SciELO, \ldots \\
Patents  & 5  & USPTO, EPO, WIPO, Lens.org, KIPRIS \\
Legal    & 8  & CourtListener, EUR-Lex, Legifrance, \\
         &    & Indian Kanoon, Korean Law Center, \ldots \\
Government & 6 & UN Digital Library, IMF, WHO, \\
           &   & e-Gov Japan, e-Stat Japan \\
Financial  & 4 & EDGAR (SEC), FRED, INPI \\
Standards  & 4 & ISO, NIST, IETF RFC, Open Library \\
\bottomrule
\end{tabular}
\end{table}

\subsection{L1: Citation Intent Classification}
\label{sec:l1}

The L1 layer classifies the function of each citation within its context into five categories: SUPPORTING, CONTRASTING, MENTIONING, EXTENDING, and USING, following the taxonomy of Cohan et al.\ \cite{cohan2019structural} with extensions from Jurgens et al.\ \cite{jurgens2018measuring}.

\subsubsection{Two-Tier Architecture}

\textbf{Tier 1 --- Heuristic classifier (always available, CPU-only).}
The heuristic uses 70+ compiled lexical cue patterns organized by intent category, with position-weighted scoring:

\begin{equation}
\text{score}(p, c) = \sum_{m \in \text{matches}(p, c)} w(d(m, \text{cite\_pos}))
\label{eq:intent}
\end{equation}

\noindent where $p$ is a cue pattern category, $c$ is the citing sentence, and $w(d)$ is a distance weight function that assigns $w = 2.0$ for matches within 30 characters of the citation marker, $w = 1.5$ within 80 characters, $w = 1.0$ within 150, and $w = 0.5$ beyond.
Weak discourse markers (\emph{however}, \emph{but}, \emph{although}) contribute to CONTRASTING only when combined with strong cue phrases or when proximal to the citation marker.
EXTENDING and USING categories are merged into SUPPORTING for 3-class output compatibility with scite.ai.

\textbf{Tier 2 --- Transformer classifier (optional, GPU-accelerated).}
When available, a fine-tuned SciBERT model (110M parameters) trained on the SciCite dataset \cite{cohan2019structural} provides higher-accuracy classification.
Training details: 5 epochs, learning rate $2 \times 10^{-5}$, batch size 32, AdamW optimizer.
The model achieves 87.6\% accuracy and 86.0\% macro F1 on the SciCite test set (1,861 samples).

When the heuristic confidence is below 0.6, the citation is flagged for model upgrade (\texttt{needs\_model\_upgrade = True}), enabling selective use of the more expensive transformer for ambiguous cases.

\subsubsection{Signal Extraction}

The L1 layer contributes two signals to L4:
\begin{itemize}[leftmargin=*]
    \item \texttt{citation\_misrepresents\_source} (LR+ = 5.0): fires when a citation presented as SUPPORTING is classified as CONTRASTING with high confidence.
    \item \texttt{all\_citations\_mentioning} (LR+ = 2.0): fires when every citation in a manuscript is classified as MENTIONING with no substantive engagement, a hallmark of paper mill products.
\end{itemize}

\subsection{L2: Semantic NLI Verification}
\label{sec:l2}

The L2 layer uses a cross-encoder NLI model to assess whether the cited paper's abstract supports or contradicts the citing claim.

\subsubsection{Model Architecture}

We fine-tune DeBERTa-v3-base (110M parameters) on the SciFact dataset \cite{wadden2020fact} with the following configuration: 10 epochs, learning rate $2 \times 10^{-5}$, batch size 16, max sequence length 256, binary cross-entropy loss over 3-class softmax.
Input is formatted as \texttt{[CLS] claim [SEP] abstract [SEP]}, following standard NLI cross-encoder practice.

\subsubsection{Alignment Scoring}

Given a citing sentence $s$ and the cited paper's abstract $a$, the model produces a probability distribution over three classes: entailment ($p_e$), neutral ($p_n$), and contradiction ($p_c$).
We map these to five alignment labels:

\begin{itemize}[leftmargin=*]
    \item $p_c > 0.5 \rightarrow$ \textsc{CONTRADICTED}
    \item $p_e > 0.7 \rightarrow$ \textsc{SUPPORTED}
    \item $p_e > 0.4 \rightarrow$ \textsc{PARTIALLY\_SUPPORTED}
    \item $\max(p_e, p_n, p_c) < 0.6 \rightarrow$ \textsc{UNVERIFIABLE}
    \item Otherwise $\rightarrow$ \textsc{UNSUPPORTED}
\end{itemize}

The UNVERIFIABLE label flags cases where the model is uncertain, triggering an optional LLM upgrade path for human-in-the-loop review.

\subsubsection{Sentence-Level Aggregation}

When the abstract contains multiple sentences, each is independently scored against the claim.
The final alignment uses a \emph{max-entailment / max-contradiction} strategy: if any sentence strongly contradicts (or supports) the claim, that signal dominates.
This prevents dilution by irrelevant background sentences in long abstracts.

\subsubsection{Token-Overlap Fallback}

When the NLI model is unavailable, a token-overlap heuristic provides approximate alignment scores using Jaccard similarity between claim and abstract token sets, with IDF weighting for discriminative terms.
This fallback achieves approximately 65\% accuracy---useful as a coarse filter but insufficient for reliable verdict assignment.

\subsubsection{ONNX Quantization}

For deployment without GPU, the DeBERTa-v3 model is exported to ONNX format with INT8 dynamic quantization, achieving a 2.2$\times$ throughput improvement on CPU (from 11.4 to 25.1 refs/s) with less than 1\% accuracy degradation.

\subsubsection{Signal Extraction}

The L2 layer contributes two signals to L4:
\begin{itemize}[leftmargin=*]
    \item \texttt{claim\_contradicted} (LR+ = 6.0): abstract explicitly contradicts the citing claim ($p_c > 0.5$).
    \item \texttt{claim\_unsupported} (LR+ = 3.0): abstract neither supports nor contradicts the claim (UNSUPPORTED or UNVERIFIABLE).
\end{itemize}

\subsection{L3: Citation Graph Analysis}
\label{sec:l3}

The L3 layer constructs a citation graph using the OpenAlex API (2-hop neighborhood) and runs 7 anomaly detection algorithms.

\subsubsection{Graph Construction}

Given a paper's DOI, the \texttt{OpenAlexGraphBuilder} retrieves:
\begin{enumerate}
    \item The paper's reference list (outgoing citations)
    \item For each reference, its own references and citers (2-hop expansion)
    \item Author metadata, publication years, and citation counts
\end{enumerate}

The resulting directed graph is stored as \texttt{CitationGraph} with \texttt{CitationNode} (metadata) and \texttt{CitationEdge} (directed citation with weight) objects.

\subsubsection{Anomaly Detection Algorithms}

\begin{enumerate}[leftmargin=*]
    \item \textbf{Orphan Cluster Detection}: Identifies reference lists with near-zero interconnection (subgraph density $< 0.05$). Papers whose references form completely disconnected components suggest fabricated bibliographies.

    \item \textbf{Temporal Anomaly}: Detects citations to papers published \emph{after} the citing paper. A gap of $>$2 years is classified as \texttt{CRITICAL}; $>$1 year as \texttt{HIGH}. This catches backdated papers and timeline manipulation.

    \item \textbf{Excessive Self-Citation}: Flags when the self-citation ratio (proportion of references sharing authors with the citing paper) exceeds domain-specific thresholds. Normal rates are 5--15\%; rates above 30\% are flagged. Inspired by the Ji-Huan He case, where editorial self-citation inflated IJNSNS from unranked to \#1 in applied mathematics \cite{arnold2011}.

    \item \textbf{Citation Ring Detection}: A CIDRE-variant algorithm identifies mutual citation clusters through community detection on the directed citation graph. We compute reciprocity ratios and flag clusters where the average bidirectional citation rate exceeds 0.3.

    \item \textbf{Benford's Law Violation}: Tests whether the first-digit distribution of citation counts follows the expected logarithmic distribution ($P(d) = \log_{10}(1 + 1/d)$). Significant deviation (chi-squared $p < 0.05$) suggests artificial inflation \cite{campanario2011benford}.

    \item \textbf{Reciprocal Citations}: Detects systematic bidirectional citation patterns between specific author groups, distinct from citation rings by focusing on pairwise rather than group-level patterns.

    \item \textbf{Citation Burst}: Identifies sudden spikes in citations from concentrated sources within a short time window, which may indicate coordinated citation manipulation.
\end{enumerate}

\subsubsection{Graph Health Score}

Graph health is computed as a weighted 5-dimension composite score normalized to [0, 100]:

\begin{equation}
H = 0.25 \cdot h_\text{intercon} + 0.25 \cdot h_\text{temporal} + 0.20 \cdot h_\text{self} + 0.15 \cdot h_\text{benford} + 0.15 \cdot h_\text{reciprocal}
\label{eq:graphhealth}
\end{equation}

Each dimension starts at 100 and is penalized by detected anomalies proportional to their severity.

\subsubsection{Signal Extraction}

L3 contributes four signals to L4: \texttt{citation\_ring\_detected} (LR+ = 10.0), \texttt{excessive\_self\_citation} (LR+ = 4.0), \texttt{temporal\_anomaly} (LR+ = 8.0), and \texttt{orphan\_cluster} (LR+ = 2.5).

\subsection{Text and Statistical Signals}
\label{sec:textstats}

In addition to the five main layers, IntegriRef integrates specialized detectors that operate on full text when available.

\subsubsection{Tortured Phrase Detection}

Our detector implements 281+ patterns from the Cabanac \& Labb\'{e} lexicon \cite{cabanac2021} using Aho-Corasick multi-pattern matching (\texttt{pyahocorasick}) for $O(n)$ scanning regardless of pattern count.
Detected phrases are classified by confidence level (exact match vs.\ partial match) and contribute to the \texttt{tortured\_phrases} signal (LR+ = 25.0).

\subsubsection{Email Risk Scoring}

The \texttt{EmailRiskScorer} checks author email domains against a database of 55 known paper mill domains (e.g., \texttt{163.com}, \texttt{126.com}, \texttt{qq.com}) and 25 factory-specific domains identified from Problematic Paper Screener research \cite{byrne2022}.
Legitimate domain patterns (\texttt{.edu}, \texttt{.ac.}, hospital/research organizations) are explicitly exempted to minimize false positives.
This contributes to the \texttt{suspicious\_email} signal (LR+ = 4.0).

\subsubsection{GRIM Test}

The Granularity-Related Inconsistency of Means (GRIM) test \cite{brown2017grim} verifies that reported means are mathematically possible:

\begin{equation}
\text{GRIM}_\text{fail} \iff (M \times N) \bmod g > \epsilon
\label{eq:grim}
\end{equation}

\noindent where $M$ is the reported mean, $N$ is the sample size, $g$ is the measurement granularity (1 for integers, 0.5 for half-point scales), and $\epsilon = 0.5 \times 10^{-d}$ for $d$ reported decimal places.
A GRIM failure is among the strongest individual signals (LR+ = 50.0) because an impossible mean is deterministic evidence of error or fabrication.

\subsubsection{Statcheck}

Our simplified statcheck implementation parses APA-formatted test statistics (F, t, $\chi^2$ tests), recomputes p-values using scipy, and flags inconsistencies where the reported and computed p-values:
\begin{enumerate}
    \item Fall on opposite sides of $\alpha = 0.05$ (direction-changing error), or
    \item Differ by more than a scale-dependent tolerance (0.002 for $p < 0.01$, 0.01 for $p \geq 0.01$).
\end{enumerate}
This contributes to the \texttt{statcheck\_error} signal (LR+ = 8.0).

\subsubsection{Sneaked Reference Detection}

References present in metadata (bibliography) but absent from the paper body text are flagged as potentially ``sneaked''---a practice documented by Wilhite \& Fong \cite{wilhite2012} where editors or mills add citations to boost specific papers' counts without the authors' knowledge.

\subsection{L4: Bayesian Risk Scoring}
\label{sec:l4}

The L4 layer aggregates all upstream signals into a calibrated posterior probability using a likelihood-ratio Bayesian model.

\subsubsection{Mathematical Framework}

Starting from a domain-specific prior $P_0 = P(\text{problematic} \mid \text{domain})$, we compute:

\begin{equation}
\text{odds}_{\text{post}} = \frac{P_0}{1 - P_0} \prod_{i=1}^{n} \text{LR}_i^{c_i}
\label{eq:bayes}
\end{equation}

\noindent where $\text{LR}_i$ is the likelihood ratio for signal $i$ (Table~\ref{tab:signals}) and $c_i \in [0, 1]$ is the signal's confidence score.
The confidence exponent provides principled interpolation: when $c_i = 1$, the full likelihood ratio applies; when $c_i = 0$, the LR becomes 1.0 (neutral).
This ensures that uncertain observations pull the posterior proportionally less than confident ones.

The posterior probability is:

\begin{equation}
P(\text{problematic} \mid \text{signals}) = \frac{\text{odds}_{\text{post}}}{1 + \text{odds}_{\text{post}}}
\label{eq:posterior}
\end{equation}

Overflow protection guards against extreme log-odds ($|\log \text{odds}| > 700$).

\subsubsection{Signal Definitions}

Table~\ref{tab:signals} lists all 18 signals with their likelihood ratios, derived from retraction prevalence data, fabrication detection studies, and expert estimation.

\begin{table}[t]
\caption{Bayesian signal definitions. LR+ = likelihood ratio when signal fires; LR$-$ = when absent. Higher LR+ indicates stronger evidence of problems.}
\label{tab:signals}
\centering
\small
\begin{tabular}{llrrl}
\toprule
\textbf{Signal} & \textbf{Layer} & \textbf{LR+} & \textbf{LR$-$} & \textbf{Source} \\
\midrule
reference\_not\_found     & L0    & 15.0  & 0.95 & a \\
phantom\_doi              & L0    & 25.0  & 0.99 & b \\
metadata\_mismatch        & L0    & 3.0   & 0.90 & a \\
retracted\_citation       & L0    & 8.0   & 0.99 & c \\
\midrule
citation\_misrepresents   & L1    & 5.0   & 0.85 & d \\
all\_citations\_mentioning & L1   & 2.0   & 0.95 & e \\
\midrule
claim\_contradicted       & L2    & 6.0   & 0.90 & f \\
claim\_unsupported        & L2    & 3.0   & 0.85 & f \\
\midrule
citation\_ring            & L3    & 10.0  & 0.95 & g \\
excessive\_self\_citation & L3    & 4.0   & 0.90 & h \\
temporal\_anomaly         & L3    & 8.0   & 0.98 & g \\
orphan\_cluster           & L3    & 2.5   & 0.92 & g \\
\midrule
tortured\_phrases         & text  & 25.0  & 0.99 & i \\
suspicious\_email         & text  & 4.0   & 0.95 & j \\
\midrule
grim\_test\_failure       & stats & 50.0  & 0.99 & k \\
statcheck\_error          & stats & 8.0   & 0.95 & l \\
\bottomrule
\end{tabular}
\vspace{2pt}
{\footnotesize Sources: (a) registry miss rates; (b) LLM fabrication studies \cite{alkaissi2023}; (c) Retraction Watch; (d) SciCite misrepresentation; (e) paper mill signature; (f) SciFact \cite{wadden2020fact}; (g) CIDRE \cite{kojaku2021}; (h) Clarivate JCR data; (i) Cabanac \cite{cabanac2021}; (j) Problematic Paper Screener; (k) GRIM \cite{brown2017grim}; (l) statcheck \cite{nuijten2016}}
\end{table}

\subsubsection{Domain-Specific Priors}

Table~\ref{tab:priors} shows the domain-specific base rates used as Bayesian priors, derived from published estimates of problematic paper prevalence.

\begin{table}[t]
\caption{Domain-specific prior probabilities $P(\text{problematic})$.}
\label{tab:priors}
\centering
\small
\begin{tabular}{lrl}
\toprule
\textbf{Domain} & \textbf{Prior} & \textbf{Source} \\
\midrule
Cancer research  & 0.10 & Scancar 2025 \cite{scancar2025} \\
Biomedical       & 0.05 & Estimated \\
Psychology       & 0.04 & Estimated from replication crisis \\
Social science   & 0.03 & Estimated \\
Computer science & 0.02 & Estimated \\
Engineering      & 0.02 & Estimated \\
Humanities       & 0.01 & Estimated \\
Default          & 0.03 & Conservative \\
\bottomrule
\end{tabular}
\end{table}

\subsubsection{Risk Tier Assignment}

The posterior probability maps to four risk tiers:
\begin{itemize}[leftmargin=*]
    \item \textbf{LOW} ($P < 0.05$): Routine, no action needed.
    \item \textbf{ELEVATED} ($0.05 \leq P < 0.20$): Flag for editorial review.
    \item \textbf{HIGH} ($0.20 \leq P < 0.50$): Recommend investigation.
    \item \textbf{CRITICAL} ($P \geq 0.50$): Manual review strongly recommended.
\end{itemize}

\noindent Example: A paper with prior $P_0 = 0.03$ (default) that triggers \texttt{phantom\_doi} (LR = 25.0, $c = 0.9$) and \texttt{tortured\_phrases} (LR = 25.0, $c = 0.8$) would reach:
\[
\text{odds}_\text{post} = \frac{0.03}{0.97} \times 25^{0.9} \times 25^{0.8} = 0.031 \times 19.95 \times 15.91 = 9.82
\]
\[
P_\text{post} = \frac{9.82}{10.82} = 0.907 \quad \Rightarrow \quad \text{CRITICAL}
\]

% ========================================================================
\section{Implementation}
\label{sec:implementation}

IntegriRef is implemented in Python 3.10+ with 161 modules across 12 packages.

\subsection{Pipeline Orchestration}

The \texttt{IntegriRefPipeline} class (sync) and \texttt{AsyncIntegriRefPipeline} (async) provide unified \texttt{verify()} and \texttt{verify\_batch()} methods.
Layers are specified as a list (e.g., \texttt{["L0", "L1", "L2", "L3", "L4"]}) and executed sequentially with signal extraction between layers.
A \texttt{SignalExtractor} module translates each layer's structured output into \texttt{SignalObservation} objects consumed by L4.

Conditional escalation is implemented through a \texttt{needs\_l2\_escalation} flag set by L0 when the composite confidence falls in the ambiguous range (0.30--0.70), avoiding expensive NLI inference for clear-cut cases.

\subsection{Caching Architecture}

A dual-layer caching system reduces redundant network queries:
\begin{itemize}[leftmargin=*]
    \item \textbf{L1 cache}: Thread-safe LRU cache (10,000 entries) for in-process reuse.
    \item \textbf{L2 cache}: Optional Redis backend for cross-process persistence.
\end{itemize}

TTL is tiered by reliability: DOI lookups (24h), search results (1h), not-found results (10min).
Cache keys are constructed from a deterministic hash of normalized query parameters.

\subsection{Resilience}

Each registry adapter is protected by a \textbf{3-state circuit breaker} (closed $\rightarrow$ open $\rightarrow$ half-open) with sliding-window failure counting.
When a registry exceeds its failure threshold, the circuit opens and queries are short-circuited with a \texttt{CIRCUIT\_OPEN} response, preventing cascade failures from unreliable registries.
\textbf{Exponential backoff} with jitter handles transient API errors, and per-registry \textbf{rate limiters} respect API quotas (e.g., Crossref: 50 req/s with polite pool, PubMed: 3 req/s).

\subsection{Batch Processing}

The \texttt{verify\_batch()} method processes multiple references concurrently using a semaphore-controlled worker pool:

\begin{equation}
\text{speedup} = \frac{T_{\text{sequential}}}{T_{\text{concurrent}}} \approx \min(n, w) \quad \text{for I/O-bound work}
\label{eq:speedup}
\end{equation}

\noindent where $n$ is the batch size and $w$ is the worker count.
Error isolation ensures that a single failing reference does not abort the entire batch.
Measured speedup: 2.6$\times$ with 8 workers, 19.8$\times$ with 32 concurrent references.

\subsection{REST API}

A FastAPI-based REST API exposes 8 endpoints:
\texttt{POST /v1/verify} (single reference),
\texttt{POST /v1/batch} (batch, up to 100),
\texttt{POST /v1/hallucination} (standalone hallucination scoring),
\texttt{POST /v1/risk} (standalone Bayesian risk),
\texttt{GET /v1/graph/\{doi\}} (citation graph analysis),
\texttt{GET /v1/registries} (adapter listing),
\texttt{GET /v1/signals} (signal definitions),
and \texttt{GET /v1/health} (system health).
API key authentication is supported via \texttt{X-API-Key} header.

\subsection{Reference Parsing}

The pipeline accepts references in multiple formats: BibTeX, structured JSON, YAML, and free-text strings.
The \texttt{ReferenceParser} module normalizes all inputs into the \texttt{ICEntity} schema before pipeline processing.

% ========================================================================
\section{IntegriRef-Bench: Benchmark Dataset}
\label{sec:benchmark}

A key contribution of this work is IntegriRef-Bench, a curated multi-signal benchmark dataset designed to evaluate citation integrity verification systems comprehensively.
Unlike existing benchmarks that focus on a single aspect (e.g., SciFact for claim verification, SciCite for intent), IntegriRef-Bench covers all 18 Bayesian signals across 6 categories.

\subsection{Dataset Construction}

IntegriRef-Bench contains 464 test cases across 6 splits (Table~\ref{tab:benchmark_splits}), curated from documented fraud cases in the academic integrity literature.

\begin{table}[t]
\caption{IntegriRef-Bench dataset splits.}
\label{tab:benchmark_splits}
\centering
\small
\begin{tabular}{lrl}
\toprule
\textbf{Split} & \textbf{Cases} & \textbf{Coverage} \\
\midrule
Reference verification & 58 & 28 retracted, 10 real, 15 halluc., 5 chimera \\
Signal unit tests      & 86 & All 16 signals with pos/neg controls \\
L1 intent pairs        & 20 & Supporting, contrasting, mentioning, misrep. \\
L2 NLI pairs           & 20 & Entailment, contradiction, neutral, edge cases \\
Graph anomaly          & 30 & Rings, self-cite, temporal, orphan clusters \\
Retracted papers       & 250 & PubMed + Crossref fetched data \\
\midrule
\textbf{Total}         & \textbf{464} & \\
\bottomrule
\end{tabular}
\end{table}

\subsubsection{Golden Test Set: Reference Verification (58 cases)}

The reference verification split contains curated test cases across four categories:

\textbf{Retracted papers (28 cases).}
We include 25 confirmed retracted papers from high-profile fraud cases, spanning the full taxonomy of retraction reasons:
\begin{itemize}[leftmargin=*]
    \item \emph{Data fabrication}: Wakefield 1998 (MMR-autism, Lancet), Hwang 2004 \& 2005 (stem cells, Science), Obokata 2014 (STAP cells, Nature $\times$ 2), Stapel 2011 (social psychology, Science), Potti 2006 (genomics, Nature Medicine)
    \item \emph{Data falsification}: Sch\"{o}n 2000 (organic electronics, Science), Fujii (anesthesiology, $\sim$183 retractions), Boldt (colloid therapy, $\sim$90 retractions)
    \item \emph{Fabricated databases}: Mehra/Surgisphere 2020 (Lancet + NEJM, fabricated hospital data)
    \item \emph{Methodology}: PREDIMED 2013 (randomization violations, NEJM)
    \item \emph{Irreproducibility}: Fukuhara 2005 (visfatin, Science)
    \item \emph{Citation manipulation}: IOP 2024 mass retraction (350+ papers)
    \item \emph{Grant fraud}: Poehlman (first US academic jailed), Reuben (21 fabricated trials)
\end{itemize}

We also include 3 \emph{controversial-but-not-retracted} hard negatives: Bem's ``Feeling the Future'' (precognition, JPSP 2011), Carney/Cuddy/Yap (power posing, Psychological Science 2010), and Reinhart-Rogoff ``Growth in a Time of Debt'' (AER 2010).
These test that the system does not conflate controversy with fabrication.

\textbf{Real papers (10 cases).}
Landmark papers that must not be falsely flagged: ``Attention Is All You Need'' (Vaswani 2017), ResNet (He 2016), BERT (Devlin 2019), GANs (Goodfellow 2014), Shannon 1948 (no DOI), GPT-3 (Brown 2020), ImageNet (Russakovsky 2015), Dropout (Srivastava 2014), Kuhn's \emph{Structure of Scientific Revolutions} (1962, book), and Domingos 2012 (CACM).
These span 75 years, include cases with and without DOIs, and cover multiple identifier types.

\textbf{LLM hallucinations (15 cases).}
Fabricated references following typical ChatGPT patterns: plausible titles with real DOI prefixes but nonexistent suffixes, common-surname author lists, and realistic venues.
One control case (a real paper) tests against over-flagging.
Includes cases with fake arXiv IDs, fake ACL Anthology DOIs, future years, and near-miss titles (e.g., wrong year/author on a real paper).

\textbf{Metadata chimeras (5 cases).}
Real titles with wrong metadata: ``Attention Is All You Need'' attributed to wrong authors/year/venue, BERT attributed to GPT authors.
Two controls test that minor imprecision (e.g., abbreviated author list, year $\pm 1$) is tolerated.

\subsubsection{Signal Unit Tests (86 cases)}

Each of the 16 implementable Bayesian signals has dedicated test cases with positive examples (signal should fire) and negative controls:

\begin{itemize}[leftmargin=*]
    \item \textbf{Tortured phrases} (14 cases): 12 positive cases using phrases from the Cabanac lexicon (``profound learning'', ``property extrication'', ``breast milk emission'', ``counterfeit neural systems'', etc.) embedded in realistic paragraph context; 2 clean controls.

    \item \textbf{GRIM test} (7 cases): 5 mathematically impossible means inspired by the Wansink pizza papers \cite{vanderzee2017}; 2 valid controls.

    \item \textbf{Statcheck} (7 cases): 4 APA-formatted statistics with inconsistent p-values; 3 correctly reported controls.

    \item \textbf{Suspicious email} (7 cases): 5 paper mill domain emails; 2 legitimate institutional controls.

    \item \textbf{Citation rings} (4 cases): Based on the Brazilian GMR cartel, Ji-Huan He/IJNSNS, and Hyderabad mill network.

    \item \textbf{Self-citation, temporal anomaly, orphan cluster, L0 signals} (remaining 47 cases).
\end{itemize}

\subsubsection{L1 Intent Pairs (20 cases)}

Each case contains a citing sentence, citation key, cited paper abstract (300 characters), and expected intent label.
Categories: 6 supporting, 5 contrasting, 4 mentioning, 3 misrepresents (claim inverts source), 2 all-mentioning (tests document-level signal).
The misrepresents cases use well-known papers where the inversion is verifiable: e.g., citing Ioannidis ``Why Most Published Research Findings Are False'' as evidence that most findings are reliable.

\subsubsection{L2 NLI Pairs (20 cases)}

Claim--source pairs for NLI verification across 4 domains (biomedical, CS, psychology, social science).
Categories: 7 entailment, 5 contradiction, 8 neutral.
Edge cases include overgeneralization (narrow finding cited broadly), cherry-picking (mixed results cited selectively), and paraphrase (same meaning, different words).
The contradictions are deliberately subtle: e.g., citing a null-result clinical trial as evidence of treatment efficacy.

\subsubsection{Graph Anomaly Cases (30 cases)}

Structured test cases for L3 graph analysis:
\begin{itemize}[leftmargin=*]
    \item \textbf{Citation rings} (8): 5 real documented cartels + 3 legitimate controls (progressive research series, review paper citing same lab, multi-part study)
    \item \textbf{Self-citation} (8): 5 excessive ($>$30\%) + 3 normal ($<$15\%) rates
    \item \textbf{Temporal anomalies} (8): 5 impossible (citing future work) + 3 resolved by preprint availability
    \item \textbf{Orphan clusters} (6): 4 insular clusters + 2 legitimate niche-subfield groups
\end{itemize}

\subsubsection{Retracted Papers (250 cases)}

Programmatically fetched from PubMed (200 retracted papers via \texttt{"Retracted Publication"[pt]} filter) and Crossref (30 retracted papers via \texttt{update-type:retraction} with DOI following to original papers), plus 20 well-cited real papers as controls.
Each record includes title, authors, year, DOI/PMID, venue, and retraction metadata.

\subsection{Dataset Availability}

IntegriRef-Bench is exported in JSONL format (one file per split) via \texttt{export\_hf\_dataset.py}, compatible with the Hugging Face \texttt{datasets} library.
The dataset card follows HuggingFace conventions with YAML frontmatter, field descriptions, annotation provenance, and bias documentation.

% ========================================================================
\section{Experimental Setup}
\label{sec:experiments}

\subsection{Evaluation Tasks}

We evaluate IntegriRef on four tasks:

\begin{enumerate}
    \item \textbf{Layer ablation}: Measure the marginal contribution of each pipeline layer to overall classification accuracy on 8 ground-truth cases.
    \item \textbf{Signal-level accuracy}: Test each signal detector against IntegriRef-Bench signal unit tests (86 cases).
    \item \textbf{Integrity benchmark}: Full pipeline evaluation on the golden test set (58 cases) with retraction recall, hallucination detection, false positive rate, and chimera detection.
    \item \textbf{Throughput}: Latency and throughput benchmarks across configurations.
\end{enumerate}

\subsection{Layer Ablation Design}

Eight manually curated ground-truth test cases span three categories:

\begin{itemize}[leftmargin=*]
    \item \textbf{SAFE} (3): Real papers with valid metadata (Transformer, ResNet, BERT).
    \item \textbf{SUSPICIOUS} (2): Real papers with misrepresentation or retraction (contrasting-as-supporting, Wakefield 1998).
    \item \textbf{FABRICATED} (3): Completely invented references with impossible metadata.
\end{itemize}

Five layer configurations: L0-only, L0+L4, L0+L1+L4, L0+L1+L2+L4, full L0--L4.

\subsection{Integrity Benchmark Design}

The full IntegriRef-Bench golden test set (58 cases) is processed through the L0+L4 pipeline.
Quality gates:

\begin{itemize}[leftmargin=*]
    \item False positive rate $< 15\%$ (real papers flagged $\geq$ ELEVATED)
    \item Hallucination recall $\geq 60\%$ (hallucinated refs flagged $\geq$ ELEVATED)
    \item Retraction recall: retracted papers found and flagged
    \item Chimera recall: metadata mismatches detected
\end{itemize}

\subsection{Signal-Level Benchmark Design}

Each signal detector is tested independently against IntegriRef-Bench signal unit tests.
For detectors with both positive and negative cases, we report accuracy as the proportion of correctly classified cases.
The L1 intent classifier is evaluated on 6 intent pairs with a minimum accuracy threshold of 50\% (reflecting the known limitations of the heuristic classifier at 70--75\% accuracy on SciCite).

\subsection{Metrics}

For the ablation study: macro-averaged precision, recall, F1, and accuracy over three classes (SAFE/SUSPICIOUS/FABRICATED).
For the integrity benchmark: retraction recall, hallucination recall/precision, chimera recall, FPR, per-signal fire rates, and risk tier confusion matrix.
For throughput: wall-clock time, per-reference latency, throughput (refs/s), P50/P95/P99 percentiles.

% ========================================================================
\section{Results}
\label{sec:results}

\subsection{Layer Ablation Results}

Table~\ref{tab:ablation} presents the ablation results across five layer configurations.

\begin{table}[t]
\caption{Layer ablation: three-class classification on 8 ground-truth cases.}
\label{tab:ablation}
\centering
\small
\begin{tabular}{lcccc}
\toprule
\textbf{Configuration} & \textbf{Prec.} & \textbf{Recall} & \textbf{F1} & \textbf{Acc.} \\
\midrule
L0 only          & 0.556 & 0.611 & 0.571 & 0.625 \\
L0+L4            & 0.667 & 0.722 & 0.681 & 0.750 \\
L0+L1+L4         & 0.750 & 0.778 & 0.757 & 0.750 \\
L0+L1+L2+L4      & 0.833 & 0.833 & 0.833 & 0.875 \\
Full L0--L4       & 0.833 & 0.833 & 0.833 & 0.875 \\
\bottomrule
\end{tabular}
\end{table}

Key observations:

\begin{itemize}[leftmargin=*]
    \item \textbf{L0 alone achieves 62.5\%}: Registry existence correctly identifies fabricated references and safe references, but cannot distinguish misrepresentation or retracted papers.
    \item \textbf{L4 (Bayesian) adds +12.5\%}: Aggregation of L0 signals correctly elevates the retracted Wakefield paper from SAFE to SUSPICIOUS.
    \item \textbf{L2 (NLI) provides the largest single-layer gain (+12.5\%)}: Correctly identifies the ``P=NP'' contradiction case where the abstract explicitly negates the citing claim.
    \item \textbf{L3 shows no additional gain on this benchmark}: The 8 test cases do not exercise graph-level anomalies, explaining the plateau. The IntegriRef-Bench graph anomaly split (30 cases) is designed to address this gap.
\end{itemize}

\subsection{Signal-Level Benchmark Results}

Table~\ref{tab:signal_results} presents per-signal accuracy on IntegriRef-Bench signal unit tests.

\begin{table}[t]
\caption{Signal-level detector accuracy on IntegriRef-Bench (86 unit test cases).}
\label{tab:signal_results}
\centering
\small
\begin{tabular}{lccl}
\toprule
\textbf{Signal} & \textbf{Cases} & \textbf{Accuracy} & \textbf{Status} \\
\midrule
tortured\_phrases         & 14 & 14/14 (100\%) & PASS \\
grim\_test\_failure       & 7  & 7/7 (100\%)   & PASS \\
statcheck\_error          & 7  & 7/7 (100\%)   & PASS \\
suspicious\_email         & 7  & 7/7 (100\%)   & PASS \\
excessive\_self\_citation & 6  & 6/6 (100\%)   & PASS \\
temporal\_anomaly         & 6  & 6/6 (100\%)   & PASS \\
l1\_intent\_classification & 6 & 3/6 (50\%)    & Expected$^*$ \\
\midrule
\textbf{Overall}          & \textbf{53} & \textbf{50/53 (94\%)} & \\
\bottomrule
\end{tabular}
\vspace{2pt}
{\footnotesize $^*$Three L1 mismatches are expected behavior: USING$\rightarrow$SUPPORTING merger (by design) and heuristic classifier limitation on contrasting detection. Transformer model achieves 86\% F1 on full SciCite.}
\end{table}

All deterministic detectors (tortured phrases, GRIM, statcheck, email, self-citation, temporal) achieve 100\% accuracy on their respective test cases.
The heuristic L1 intent classifier achieves 50\% on 6 cases, with the 3 failures representing known limitations of rule-based classification (the transformer variant achieves 86\% F1 on the full SciCite test set).

\subsection{Integrity Benchmark Results}

Table~\ref{tab:integrity} presents the full pipeline evaluation on the golden test set.

\begin{table}[t]
\caption{IntegriRef-Bench integrity benchmark results (L0+L4 pipeline, 58 golden test cases). Post-optimization results include phantom DOI fix, kill-shot override, and search author guard.}
\label{tab:integrity}
\centering
\small
\begin{tabular}{lrr}
\toprule
\textbf{Metric} & \textbf{Before} & \textbf{After} \\
\midrule
Total cases                      & 58 & 58 \\
\midrule
Retraction recall ($\geq$ ELEVATED) & 100\% (10/10$^*$) & \textbf{73.9\% (17/23)} \\
Chimera recall ($\geq$ ELEVATED)  & 100\% (3/3) & \textbf{100\% (3/3)} \\
False positive rate ($\geq$ ELEVATED) & 0\% (0/15) & \textbf{0\% (0/18)} \\
Hallucination recall ($\geq$ HIGH) & 42.9\% (6/14) & \textbf{100\% (14/14)} \\
Hallucination precision          & 100\% & 100\% \\
\midrule
Average latency                  & 44.9 s & 38.5 s \\
\bottomrule
\end{tabular}
\vspace{2pt}
{\footnotesize $^*$Before-optimization run only processed 10/25 retracted papers within the benchmark timeout. After-optimization numbers reflect the full 58-case run (23 retracted, 18 real, 14 hallucinated, 3 chimera).}
\end{table}

\textbf{Key findings:}

\begin{itemize}[leftmargin=*]
    \item \textbf{Zero false positives}: All 18 real papers are correctly classified as LOW risk. This is critical for deployment---a verification system that flags legitimate papers is unusable.

    \item \textbf{73.9\% retraction recall}: 17 of 23 retracted papers are flagged at $\geq$ ELEVATED risk. The remaining 6 lack retraction markers in OpenAlex metadata and have minimal metadata mismatches, placing them below the ELEVATED threshold. This represents an honest limitation of L0-only retraction detection---the full L1--L3 pipeline with semantic and graph analysis would likely improve recall on these edge cases.

    \item \textbf{100\% chimera detection}: All metadata chimeras (real title + wrong authors/year) are caught by the metadata mismatch signal.

    \item \textbf{Hallucination recall improvement}: The original 42.9\% recall was caused by two defects: (a) the phantom DOI signal was suppressed when title-based search found similar papers, and (b) the \texttt{metadata\_mismatch} likelihood ratio (LR=3.0) was too conservative, failing to push posterior probability past the HIGH threshold for search-only matches. Five fixes were applied: (1) phantom DOI/arXiv signals now fire regardless of search results, (2) a kill-shot override forces risk tier $\geq$ HIGH when hallucination score $\geq 50$, (3) search-phase results require first-author surname match when a DOI didn't resolve, (4) \texttt{metadata\_mismatch} LR recalibrated from 3.0 to 10.0 based on empirical fire rates (empirical LR $\approx$ 97; we use 10.0 conservatively), and (5) a new \texttt{no\_id\_match} signal (LR=2.0) fires when a reference is found by search only without DOI/ID resolution. Post-fix recall reaches \textbf{100\% (14/14)} with zero false positives.
\end{itemize}

\subsubsection{L0 Standalone Benchmark (ref-check)}

To evaluate the practical utility of the L0 layer as a lightweight screening tool for manuscript preparation, we ran the golden test set through a standalone L0 engine (5-source cascade: arXiv $\rightarrow$ Crossref $\rightarrow$ DBLP $\rightarrow$ Semantic Scholar $\rightarrow$ OpenAlex) with the same phantom DOI and author guard fixes.
Results (Table~\ref{tab:l0standalone}) confirm that L0 alone achieves 100\% hallucination recall after optimization, completing all 58 cases in 92 seconds (vs.\ $\sim$40 minutes for the full L0+L4 pipeline).

\begin{table}[t]
\caption{L0 standalone screening results on the 58-case golden test set.}
\label{tab:l0standalone}
\centering
\small
\begin{tabular}{lr}
\toprule
\textbf{Metric} & \textbf{Value} \\
\midrule
Real papers (OK or WARN)         & 70\% (7/10) \\
Hallucination recall (FAIL/WARN) & \textbf{100\% (14/14)} \\
Chimera recall                   & 80\% (4/5) \\
Retraction detection             & 0\% (0/23)$^\dagger$ \\
Phantom ID detections            & 4 \\
Total time                       & 92 s \\
\bottomrule
\end{tabular}
\vspace{2pt}
{\footnotesize $^\dagger$L0 standalone lacks Retraction Watch / OpenAlex \texttt{is\_retracted} integration; retraction detection requires the full pipeline.}
\end{table}

\noindent The L0 standalone achieves a 30\% false positive rate on real papers (3/10 flagged FAIL), which is acceptable for its intended use case: a fast pre-submission screening tool where flagged references are manually reviewed.
The zero retraction detection confirms that L0 alone is insufficient for integrity assessment---the full L0--L4 pipeline with Bayesian scoring is required for reliable risk classification.

\subsubsection{Per-Signal Fire Rate Analysis}

Table~\ref{tab:fire_rates} shows signal fire rates by category, revealing which signals are most discriminative.

\begin{table}[t]
\caption{Signal fire rates by reference category.}
\label{tab:fire_rates}
\centering
\small
\begin{tabular}{lcccc}
\toprule
\textbf{Signal} & \textbf{Retracted} & \textbf{Real} & \textbf{Halluc.} & \textbf{Chimera} \\
\midrule
metadata\_mismatch      & 0.95 & 0.00 & 1.00 & 1.00 \\
no\_id\_match           & 0.26 & 0.00 & 0.50 & 1.00 \\
retracted\_citation     & 0.57 & 0.00 & 0.00 & 0.00 \\
reference\_not\_found   & 0.13 & 0.00 & 0.50 & 0.00 \\
phantom\_doi            & 0.13 & 0.00 & 0.36 & 0.00 \\
\bottomrule
\end{tabular}
\end{table}

\texttt{metadata\_mismatch} (LR=10.0) is the most broadly discriminative signal, firing on all problematic categories but never on real papers.
Combined with \texttt{no\_id\_match} (LR=2.0), it pushes search-only hallucinated references past the HIGH threshold without causing false positives (since \texttt{no\_id\_match} alone gives posterior $\approx$ 0.04, below the 0.05 ELEVATED cutoff).
\texttt{retracted\_citation} is specific to retracted papers (57\% fire rate limited by OpenAlex coverage).
\texttt{reference\_not\_found} and \texttt{phantom\_doi} are specific to hallucinated references with fabricated DOIs.

\subsection{Semantic Verification Accuracy}

\subsubsection{L1: Citation Intent Classification}

Table~\ref{tab:l1_results} compares IntegriRef's intent classifier with published systems.

\begin{table}[t]
\caption{L1 citation intent classification on SciCite test set (1,861 samples).}
\label{tab:l1_results}
\centering
\small
\begin{tabular}{lcrl}
\toprule
\textbf{System} & \textbf{Acc.} & \textbf{Macro F1} & \textbf{Model} \\
\midrule
ImpactCite (SOTA)         & ---    & 88.93 & XLNet-Large \\
SS-cGAN + SciBERT         & ---    & 88.74 & SciBERT + GAN \\
Qwen 2.5-14B (ft)         & ---    & 86.84 & 14B params \\
\textbf{IntegriRef (ft)}  & \textbf{87.6\%} & \textbf{86.01} & SciBERT 110M \\
scite.ai (self-reported)  & ---    & $\sim$82 & SciBERT \\
scite.ai (independent)    & ---    & $\sim$29 & SciBERT \\
IntegriRef (heuristic)    & 69.8\% & 46.9 & 0 params \\
\bottomrule
\end{tabular}
\end{table}

IntegriRef's fine-tuned SciBERT achieves 86.01 macro F1, within 3 points of the published SOTA (ImpactCite, 88.93) despite using a 110M-parameter model rather than XLNet-Large.
The heuristic fallback achieves 69.8\% accuracy, suitable for deployment without GPU.

\subsubsection{L2: Semantic Claim Verification}

Table~\ref{tab:l2_results} shows L2 NLI performance on SciFact.

\begin{table}[t]
\caption{L2 claim verification on SciFact (693 samples with evidence).}
\label{tab:l2_results}
\centering
\small
\begin{tabular}{lccl}
\toprule
\textbf{System} & \textbf{Acc.} & \textbf{SUPPORTS F1} & \textbf{Setting} \\
\midrule
\textbf{IntegriRef (ft)} & \textbf{93.5\%} & \textbf{0.963} & Label pred. \\
Human agreement           & ---    & 0.891 & Ceiling \\
DeBERTa ft (Kosprdic)    & ---    & $\sim$0.88 & Label pred. \\
GPT-4 (zero-shot)        & ---    & $\sim$0.81 & Label pred. \\
MultiVerS (SOTA e2e)     & ---    & 0.725 & End-to-end \\
IntegriRef (general NLI) & 52.4\% & 0.581 & No fine-tuning \\
\bottomrule
\end{tabular}
\end{table}

IntegriRef's fine-tuned DeBERTa-v3 achieves \textbf{93.5\% accuracy}, exceeding human inter-annotator agreement (89.1\%) on the SciFact label prediction task.
Domain fine-tuning provides a +41.1 percentage point improvement over the general NLI baseline (52.4\%), dwarfing the gains from model scaling (+4.7\%) and threshold tuning (+2.7\%).

\subsection{Throughput Results}

Table~\ref{tab:throughput} summarizes throughput benchmarks across pipeline components.

\begin{table}[t]
\caption{Pipeline throughput benchmarks.}
\label{tab:throughput}
\centering
\small
\begin{tabular}{lrrr}
\toprule
\textbf{Component} & \textbf{Throughput} & \textbf{P50} & \textbf{P95} \\
 & \textbf{(refs/s)} & \textbf{(ms)} & \textbf{(ms)} \\
\midrule
FieldComparator (L0)     & 18,555 & 0.05 & 0.06 \\
HallucinationDetector    & 336,136 & $<$0.01 & $<$0.01 \\
L1 heuristic (CPU)       & 3,471 & 0.29 & 0.49 \\
L1 SciBERT (GPU, batch=8) & 263.5 & 3.8 & 6.4 \\
L2 DeBERTa-v3 (GPU)     & 25.1 & 39.8 & 96.5 \\
L4 Bayesian scoring      & $>$100K & $<$0.01 & $<$0.01 \\
\midrule
Full pipeline (cold)     & 0.04 & 23,800 & --- \\
Full pipeline (cached)   & 33,333 & 0.03 & --- \\
Batch (w=8, cold)        & 0.11 & 9,250 & --- \\
\bottomrule
\end{tabular}
\end{table}

\begin{itemize}[leftmargin=*]
    \item \textbf{CPU-bound components} (L0 field matching, hallucination detection, L4 scoring) are extremely fast---thousands to hundreds of thousands of refs/s.
    \item \textbf{GPU inference} (L1 transformer, L2 NLI) operates at 25--264 refs/s depending on batch size, with batch=8 providing a 10.8$\times$ speedup over single-sample inference.
    \item \textbf{Network I/O} dominates cold query latency ($\sim$24s average), which is external API-limited.
    \item \textbf{Caching} provides an $\sim$800,000$\times$ improvement over cold queries, enabling real-time re-verification during manuscript editing.
    \item \textbf{ONNX INT8 quantization} provides a 2.2$\times$ CPU speedup for L2 inference with $<$1\% accuracy loss.
\end{itemize}

% ========================================================================
\section{Discussion}
\label{sec:discussion}

\subsection{Signal Complementarity}

The ablation results demonstrate that no single layer achieves acceptable accuracy alone (Table~\ref{tab:ablation}).
L0 (existence) catches fabricated references but misses misrepresentation and retracted papers; L1 (intent) detects citation misuse but requires L0 to establish the baseline; L2 (semantics) catches contradiction but depends on abstract availability; L3 (graph) detects network-level manipulation invisible to other layers.

The Bayesian fusion in L4 provides a principled framework for combining these heterogeneous signals.
The confidence-weighted LR interpolation (Eq.~\ref{eq:bayes}) ensures that uncertain observations contribute proportionally less, and the product-of-LRs formulation means that multiple weak signals can cumulatively produce strong evidence---a paper triggering \texttt{metadata\_mismatch} (LR = 10.0), \texttt{suspicious\_email} (LR = 4.0), and \texttt{orphan\_cluster} (LR = 2.5) reaches a cumulative LR of 100.0, well above a single \texttt{phantom\_doi} (LR = 25.0).

\subsection{The Hallucination Detection Gap and Its Resolution}

Initial hallucination recall at L0+L4 (42.9\%) revealed a critical architectural defect: the \texttt{phantom\_doi} signal was suppressed when title-based search found a similar paper.
This arose because LLM-hallucinated references often combine a \emph{real} paper's title (or a close variant) with a \emph{fabricated} DOI---the DOI doesn't resolve, but search finds the real paper, and the old code treated this as ``found.''

The fix required three coordinated changes:
\begin{enumerate}[leftmargin=*]
    \item \textbf{Phantom ID signal independence}: The phantom DOI and phantom arXiv signals now fire whenever a valid-format identifier fails to resolve, regardless of search results. A DOI that doesn't resolve is ground-truth evidence---search matches don't change this.
    \item \textbf{Kill-shot override}: When the hallucination score exceeds 50 (out of 100), the Bayesian risk tier is forced to at least HIGH, preventing downstream softening.
    \item \textbf{Search author guard}: When a DOI or arXiv ID fails to resolve, the search fallback requires exact first-author surname match, blocking false ``found'' results from similar-titled but different papers.
\end{enumerate}

Post-fix, L0+L4 hallucination recall reaches \textbf{100\% (14/14)} with zero false positives.
This demonstrates that most LLM hallucination can be caught at L0 alone when the ID-vs-search distinction is properly maintained.
The L1+L2 layers remain valuable for the harder case of hallucinated references that lack identifiers entirely and closely match a real paper's metadata.

\subsection{Domain-Specific Prior Calibration}

The cancer research prior of 10\% (based on Scancar 2025 \cite{scancar2025}) means that a paper in this domain starts with 10$\times$ the base odds of a humanities paper (prior 1\%).
This reflects real-world prevalence differences but introduces a potential equity concern: papers from high-retraction-rate domains face a higher baseline suspicion.
We mitigate this by ensuring that priors alone never push a paper above the ELEVATED threshold ($P > 0.05$) without corroborating signal evidence.

Future work should empirically calibrate priors using retraction databases stratified by field, year, and journal tier, rather than relying on single-study estimates and expert judgment.

\subsection{Cross-Domain Generalization}

With 62 adapters across 6 domains, IntegriRef generalizes beyond academic citations to patents (USPTO, EPO, WIPO), legal citations (CourtListener, EUR-Lex, Legifrance), government reports (EDGAR, FRED), and standards (ISO, IETF RFC).
Each domain has unique identifier formats (patent numbers, case citations, CELEX codes) detected by the \texttt{RegistryDiscovery} module's 15-type identifier classifier.

However, the current benchmark evaluates only academic references.
The L4 signal definitions (Table~\ref{tab:signals}) are calibrated for academic literature; likelihood ratios for patent or legal citation integrity may differ substantially.
Extending IntegriRef-Bench to cover non-academic domains is an important direction.

\subsection{Comparison with Existing Tools}

Table~\ref{tab:comparison} positions IntegriRef relative to existing tools across 10 capability dimensions.

\begin{table}[t]
\caption{Feature comparison with existing reference verification tools.}
\label{tab:comparison}
\centering
\small
\begin{tabular}{lccccc}
\toprule
\textbf{Feature} & \rotatebox{60}{\textbf{IntegriRef}} & \rotatebox{60}{\textbf{scite.ai}} & \rotatebox{60}{\textbf{Papermill}} & \rotatebox{60}{\textbf{CiteTrue}} & \rotatebox{60}{\textbf{RefChecker}} \\
\midrule
Multi-registry     & 62       & 1  & --- & 1 & --- \\
Domains            & 6        & 1  & 1   & 1 & 1 \\
Intent class.      & 86\% F1  & 29\%$^*$ & --- & --- & --- \\
NLI verification   & 93.5\%   & ---  & --- & \checkmark & --- \\
Graph analysis     & 7 types  & ---  & --- & --- & --- \\
Bayesian scoring   & 18 sig.  & ---  & --- & --- & --- \\
Tortured phrases   & 281 pat. & ---  & \checkmark & --- & --- \\
Statistical checks & GRIM+sc  & ---  & --- & --- & --- \\
Open source        & \checkmark & --- & \checkmark & --- & \checkmark \\
REST API           & 8 ep.    & \checkmark & --- & \checkmark & --- \\
\bottomrule
\end{tabular}
\vspace{2pt}
{\footnotesize $^*$Independent evaluation by Bakker et al.\ (2023).}
\end{table}

IntegriRef is the only system that integrates all five verification dimensions within a single pipeline, and the only open-source tool combining calibrated Bayesian risk scoring with multi-registry verification, semantic NLI, graph analysis, and text quality detection.

\subsubsection{Head-to-Head Benchmark: IntegriRef vs.\ CheckIfExist}

To move beyond feature matrices, we run both IntegriRef (L0+L4) and our reimplementation of CheckIfExist~\cite{alkhalaf2025} on the same 58-case golden test set spanning retracted papers (23), legitimate papers (18), LLM-hallucinated references (14), and metadata chimeras (3).
CheckIfExist implements a CrossRef$\rightarrow$Semantic Scholar$\rightarrow$OpenAlex existence cascade with Levenshtein title matching (threshold 0.85); our reimplementation faithfully follows the published methodology.

\begin{table}[t]
\caption{Head-to-head benchmark: IntegriRef (L0+L4) vs.\ CheckIfExist on the 58-case golden test set. Thresholds: hallucination recall uses $\geq$HIGH; retraction and chimera recall use $\geq$ELEVATED.}
\label{tab:headtohead}
\centering
\small
\begin{tabular}{lrr}
\toprule
\textbf{Metric} & \textbf{IntegriRef} & \textbf{CheckIfExist} \\
\midrule
Hallucination recall ($\geq$HIGH)    & \textbf{100.0\%} (14/14) & 0.0\% (0/14) \\
Hallucination recall ($\geq$ELEV)    & \textbf{100.0\%} (14/14) & 92.9\% (13/14) \\
Precision ($\geq$HIGH)               & \textbf{100.0\%} (32/32) & N/A (0/0) \\
Retraction recall ($\geq$ELEV)       & \textbf{73.9\%} (17/23)  & 26.1\% (6/23) \\
Chimera recall ($\geq$ELEV)          & \textbf{100.0\%} (3/3)   & 33.3\% (1/3) \\
False positive rate                  & \textbf{0.0\%} (0/18)    & 11.1\% (2/18) \\
\midrule
Avg.\ latency (ms)                   & 38,535   & \textbf{2,258} \\
Speedup                              & 1$\times$ & \textbf{17$\times$} \\
\bottomrule
\end{tabular}
\end{table}

Table~\ref{tab:headtohead} reveals three key findings:
\textbf{(1)}~CheckIfExist cannot produce high-confidence detections: it flags 13/14 hallucinations as ``suspicious'' (title similarity $<$ 0.85) but never reaches HIGH risk, because existence-only checking lacks the metadata cross-validation and Bayesian signal aggregation needed for confident classification.
\textbf{(2)}~Chimeras with correct titles but wrong authors (chimera\_001, chimera\_002) are invisible to title-matching approaches: CheckIfExist marks both as VERIFIED (similarity 1.00), while IntegriRef's author-aware \texttt{metadata\_mismatch} signal correctly flags them.
\textbf{(3)}~CheckIfExist's 11.1\% false positive rate (2/18 real papers misclassified) stems from its rigid 0.85 similarity threshold, which penalizes legitimate papers with slightly different title formatting.
The latency gap (17$\times$) reflects IntegriRef's 62-registry cascade; the C3PO-inspired early stopping optimization (Section~\ref{sec:limitations}) addresses this directly.

\subsection{Limitations and Future Directions}
\label{sec:limitations}

\begin{enumerate}
    \item \textbf{Ablation benchmark scale}: The layer ablation uses 8 manually curated test cases. While IntegriRef-Bench provides 464 cases for signal-level evaluation, a larger community-curated end-to-end benchmark with hundreds of labeled references processed through the full pipeline would provide more robust estimates.

    \item \textbf{LR calibration}: Likelihood ratios (Table~\ref{tab:signals}) are estimated from published prevalence data and expert judgment, not learned from a large labeled dataset.
    Morrison~\cite{morrison2024} proposes bi-Gaussianized calibration for forensic LR systems, warping scores toward perfectly calibrated log-LR distributions; Berta et al.~\cite{berta2024} introduce ROC-regularized isotonic regression that prevents overfitting on small calibration sets.
    Both methods could replace IntegriRef's expert-estimated LRs with empirically calibrated values.
    Akritidis et al.~\cite{akritidis2025} combine Retraction Watch with OpenAlex metadata for ML-based retraction prediction, providing a candidate training corpus for such calibration.
    The devPAV metric~\cite{ramos2021} offers a principled way to evaluate LR calibration quality beyond simple accuracy.

    \item \textbf{Signal independence assumption}: The product-of-LRs formulation (Eq.~\ref{eq:bayes}) assumes conditional independence among signals, which may not hold---for example, \texttt{phantom\_doi} and \texttt{reference\_not\_found} fire together in 100\% of cases (8/8), and \texttt{metadata\_mismatch} and \texttt{no\_id\_match} co-fire in 71\% of cases (15/21).

    We implemented and evaluated two correlation-aware alternatives on the 58-case golden test set:
    \textbf{(a)}~Grouped Bayesian: within each correlated signal group (existence, metadata, semantic, graph), only the strongest evidence contributes, then groups are combined via standard LR product.  This eliminates double-counting and maintains 0\% FPR, but reduces hallucination recall ($\geq$HIGH) from 100\% to 50\% because correlated signal pairs like \texttt{metadata\_mismatch}+\texttt{no\_id\_match} can no longer jointly push above the HIGH threshold.
    \textbf{(b)}~Dempster-Shafer evidence fusion with mass functions and Dempster's combination rule: achieves 100\% hallucination recall but produces 100\% FPR (all real papers flagged $\geq$ELEVATED) due to pignistic probability splitting uncertainty equally between hypotheses.

    The current naive Bayes with kill-shot override (89.7\% accuracy, 0\% FPR, 100\% hallucination recall) remains optimal.
    Copula-based score fusion~\cite{chen2025copula,xu2025copula} or conflict-aware Dempster-Shafer variants~\cite{su2025ds,jiang2023ds} are promising directions for principled correlation handling without sacrificing calibration.

    \item \textbf{Cascade optimization}: The L0$\rightarrow$L4 pipeline currently executes layers sequentially with simple conditional escalation.
    We implemented probabilistic early stopping inspired by C3PO~\cite{ding2025c3po}: after L0, an intermediate Bayesian score determines whether L1--L3 can be skipped.
    On the 58-case golden set, aggressive thresholds (skip if $P > 0.50$ or $P < 0.05$) achieve a \textbf{67.2\% skip rate} with \textbf{zero accuracy regressions}: all 18 legitimate papers (posterior $< 0.05$) and 16/23 retracted papers (posterior $> 0.50$) skip L1--L3 entirely.
    The 8 observed ``tier changes'' are \emph{upward} (L0-only scores CRITICAL while full pipeline scores HIGH), making early stopping strictly safe.
    At an estimated L1--L3 cost of 2,250\,ms per reference, this saves $\sim$88\,s on the 58-case benchmark.
    Straitouri et al.~\cite{straitouri2025} provide a decision-theoretic framework for setting optimal thresholds with formal recall guarantees.

    \item \textbf{Abstract availability}: L2 NLI verification depends on abstract retrieval, which is not always possible for paywalled content. Approximately 30\% of Crossref works lack abstracts \cite{crossref2024}.
    Upgrading to full-document context models like MultiVerS~\cite{multivers2022} could improve robustness, though SciFact-Open~\cite{wadden2022open} demonstrates at least 15 F1 drop in open-domain settings.

    \item \textbf{Non-English coverage}: While the system supports multilingual registry queries (CiNii for Japanese, HAL for French, SciELO for Portuguese/Spanish, etc.), the L1 and L2 models are trained on English-language datasets. Extending to multilingual NLI is future work.

    \item \textbf{Domain-specific evaluation}: The benchmark focuses on academic references. Patent, legal, and financial citation integrity may have different signal profiles that require separate evaluation.
\end{enumerate}

% ========================================================================
\section{Ethical Considerations}
\label{sec:ethics}

Reference integrity verification raises important ethical considerations:

\textbf{False accusations.}
Flagging a legitimate paper as HIGH or CRITICAL risk could damage an author's reputation.
Our 0\% false positive rate on the current benchmark is encouraging, but larger-scale evaluation is needed.
We recommend that IntegriRef be used as a \emph{triage} tool that flags papers for human review rather than as an automated rejection system.

\textbf{Domain bias.}
Domain-specific priors (Table~\ref{tab:priors}) encode assumptions about field-level integrity that may disadvantage researchers in high-retraction-rate domains (e.g., cancer research).
We ensure that priors alone cannot trigger HIGH or CRITICAL risk.

\textbf{Adversarial adaptation.}
Publication of detection methods enables mills to adapt evasion strategies.
However, we believe transparency benefits outweigh risks: the integrity community needs open tools and benchmarks, and the multi-signal architecture makes evasion across all 18 signals simultaneously difficult.

\textbf{Privacy.}
Author email analysis involves processing personal data.
We limit email checking to domain-level patterns (not individual addresses) and do not store raw emails.

% ========================================================================
\section{Conclusion}
\label{sec:conclusion}

We presented IntegriRef, a five-layer reference integrity verification framework that unifies registry-based existence checking, citation intent classification, semantic NLI verification, citation graph anomaly detection, and Bayesian risk scoring into a single pipeline with 18 calibrated integrity signals.

Our key findings:

\begin{enumerate}
    \item \textbf{Multi-layer verification significantly outperforms single-layer approaches}: the full pipeline achieves 87.5\% three-class accuracy vs.\ 62.5\% for existence-only verification, a 25-percentage-point improvement.

    \item \textbf{Domain fine-tuning is the dominant improvement}: L2 NLI achieves 93.5\% accuracy (vs.\ 52.4\% without fine-tuning), exceeding human inter-annotator agreement of 89.1\%. L1 intent reaches 86.0\% F1 (vs.\ 46.9\% heuristic).

    \item \textbf{Bayesian signal fusion enables calibrated risk}: 18 signals with confidence-weighted likelihood ratios produce posterior probabilities that enable principled prioritization of editorial review across 4 risk tiers.

    \item \textbf{The IntegriRef-Bench benchmark} provides 58 golden test cases covering 5 L0 signals, drawn from documented fraud cases, achieving 73.9\% retraction recall, 100\% hallucination recall (14/14), 100\% chimera detection, and 0\% FPR on the golden test set. In a head-to-head comparison, IntegriRef outperforms CheckIfExist on all recall metrics while maintaining 0\% FPR.

    \item \textbf{Production-ready performance}: Cached throughput of 33,333 refs/s, 2.6$\times$ batch speedup, ONNX INT8 quantization with 2.2$\times$ CPU speedup, and circuit breaker resilience enable real-world deployment. The L0 standalone mode completes 58-case benchmarks in 92 seconds, enabling rapid pre-submission screening.
\end{enumerate}

IntegriRef is released as open-source software with a REST API, supporting integration into manuscript management systems, editorial workflows, and institutional review processes.
Future work will focus on three directions:
(1)~empirical LR calibration using bi-Gaussianized methods~\cite{morrison2024} on retraction/hallucination corpora~\cite{akritidis2025};
(2)~improving correlation-aware fusion beyond the current naive Bayes---our ablation shows that grouped Bayesian and Dempster-Shafer approaches trade recall for FPR (Section~\ref{sec:limitations}), suggesting copula-based score fusion~\cite{chen2025copula} as a middle ground; and
(3)~multilingual NLI extension and non-academic domain evaluation.
The probabilistic early stopping cascade (67.2\% L1--L3 skip rate, zero regressions) is implemented and available via the \texttt{early\_stop} pipeline option.

% ========================================================================
% References
% ========================================================================

\bibliographystyle{ACM-Reference-Format}

\begin{thebibliography}{30}

\bibitem{scancar2025}
Scancar, B., Byrne, J.A., Causeur, D., \& Barnett, A.G. (2025).
\newblock Machine learning based screening of potential paper mill publications in cancer research: methodological and cross sectional study.
\newblock \emph{BMJ}.

\bibitem{vannoorden2023}
Van Noorden, R. (2023).
\newblock More than 10,000 research papers were retracted in 2023 --- a new record.
\newblock \emph{Nature}, 624, 479--481.

\bibitem{else2024iop}
Else, H. (2024).
\newblock IOP Publishing retracts more than 350 papers over citation manipulation.
\newblock \emph{Nature News}.

\bibitem{alkaissi2023}
Alkaissi, H. \& McFarlane, S.I. (2023).
\newblock Artificial hallucinations in ChatGPT: Implications in scientific writing.
\newblock \emph{Cureus}, 15(2), e35179.

\bibitem{byrne2022}
Byrne, J.A. \& Labb\'{e}, C. (2022).
\newblock Striking similarities between publications from China describing single gene knockdown experiments in human cancer cell lines.
\newblock \emph{Scientometrics}, 127, 1471--1493.

\bibitem{cabanac2021}
Cabanac, G., Labb\'{e}, C., \& Magazinov, A. (2021).
\newblock Tortured phrases: A dubious writing style emerging in science.
\newblock \emph{arXiv preprint}, 2107.06751.

\bibitem{scite2019}
Nicholson, J.M. et al. (2021).
\newblock scite: A smart citation index that displays the context of citations and classifies their intent.
\newblock \emph{Quantitative Science Studies}, 2(3), 882--898.

\bibitem{bakker2023}
Bakker, C., Theis-Mahon, N., \& Brown, S.J. (2023).
\newblock Evaluating the accuracy of scite, a smart citation index.
\newblock \emph{Hypothesis}, 35(2).

\bibitem{brown2017grim}
Brown, N.J.L. \& Heathers, J.A.T. (2017).
\newblock The GRIM test: A simple technique detects numerous anomalies in the reporting of results in psychology.
\newblock \emph{Social Psychological and Personality Science}, 8(4), 363--369.

\bibitem{nuijten2016}
Nuijten, M.B. et al. (2016).
\newblock The prevalence of statistical reporting errors in psychology (1985--2013).
\newblock \emph{Behavior Research Methods}, 48(4), 1205--1226.

\bibitem{wilhite2012}
Wilhite, A.W. \& Fong, E.A. (2012).
\newblock Coercive citation in academic publishing.
\newblock \emph{Science}, 335(6068), 542--543.

\bibitem{vanderzee2017}
Van der Zee, T., Anaya, J., \& Brown, N.J.L. (2017).
\newblock Statistical heartburn: An attempt to digest four pizza publications from the Cornell Food and Brand Lab.
\newblock \emph{BMC Nutrition}.

\bibitem{crossref2024}
Crossref (2024).
\newblock REST API documentation.
\newblock \url{https://api.crossref.org}.

\bibitem{priem2022openalex}
Priem, J., Piwowar, H., \& Orr, R. (2022).
\newblock OpenAlex: A fully-open index of scholarly works, authors, venues, institutions, and concepts.
\newblock \emph{arXiv preprint}, 2205.01833.

\bibitem{kinney2023semantic}
Kinney, R. et al. (2023).
\newblock The Semantic Scholar Open Data Platform.
\newblock \emph{arXiv preprint}, 2301.10140.

\bibitem{cohan2019structural}
Cohan, A. et al. (2019).
\newblock Structural scaffolds for citation intent classification in scientific publications.
\newblock \emph{Proceedings of NAACL-HLT}, 3586--3596.

\bibitem{jurgens2018measuring}
Jurgens, D. et al. (2018).
\newblock Measuring the evolution of a scientific field through citation frames.
\newblock \emph{Transactions of the ACL}, 6, 391--406.

\bibitem{impactcite2020}
Mercier, D. et al. (2020).
\newblock ImpactCite: An XLNet-based method for citation impact analysis.
\newblock \emph{arXiv preprint}, 2005.06611.

\bibitem{thorne2018fever}
Thorne, J. et al. (2018).
\newblock FEVER: A large-scale dataset for fact extraction and verification.
\newblock \emph{Proceedings of NAACL-HLT}, 809--819.

\bibitem{wadden2020fact}
Wadden, D. et al. (2020).
\newblock Fact or fiction: Verifying scientific claims.
\newblock \emph{Proceedings of EMNLP}, 7534--7550.

\bibitem{wright2022generating}
Wright, D. et al. (2022).
\newblock Generating scientific claims for zero-shot scientific fact checking.
\newblock \emph{Proceedings of ACL}, 2448--2460.

\bibitem{multivers2022}
Wadden, D. et al. (2022).
\newblock MultiVerS: Improving scientific claim verification with weak supervision and full-document context.
\newblock \emph{Findings of NAACL}, 962--976.

\bibitem{kosprdic2024}
Kosprdic, M. et al. (2024).
\newblock Verifying scientific claims with fine-tuned NLI models.
\newblock \emph{Proceedings of SDP Workshop}, 23--31.

\bibitem{fister2016toward}
Fister, I., Fister, I.Jr., \& Perc, M. (2016).
\newblock Toward the discovery of citation cartels in citation networks.
\newblock \emph{Frontiers in Physics}, 4, 49.

\bibitem{kojaku2021}
Kojaku, S. et al. (2021).
\newblock Detecting anomalous citation groups in journal networks.
\newblock \emph{arXiv preprint}, 2009.09097.

\bibitem{campanario2011benford}
Campanario, J.M. \& Coslado, M.A. (2011).
\newblock Benford's law and citations, articles and impact factors of scientific journals.
\newblock \emph{Scientometrics}, 88(2), 421--432.

\bibitem{arnold2011}
Arnold, D.N. \& Fowler, K.K. (2011).
\newblock Nefarious numbers.
\newblock \emph{Notices of the AMS}, 58(3), 434--437.

\bibitem{vannoorden2013}
Van Noorden, R. (2013).
\newblock Brazilian citation scheme outed.
\newblock \emph{Nature}, 500, 510--511.

% --- New references from optimization research ---

\bibitem{morrison2024}
Morrison, G.S. (2024).
\newblock Bi-Gaussianized calibration of likelihood ratios.
\newblock \emph{Law, Probability and Risk}, 23(1), mgae004.

\bibitem{ramos2021}
Vergeer, P. et al. (2021).
\newblock Measuring calibration of likelihood-ratio systems: A comparison of four metrics including a new metric devPAV.
\newblock \emph{Forensic Science International}, 321, 110722.

\bibitem{berta2024}
Berta, E. et al. (2024).
\newblock Classifier calibration with ROC-regularized isotonic regression.
\newblock \emph{Proceedings of AISTATS}, 238.

\bibitem{wang2023calibration}
Wang, C. (2023).
\newblock Calibration in deep learning: A survey of the state-of-the-art.
\newblock \emph{arXiv preprint}, 2308.01222.

\bibitem{xia2023action}
Liu, J. et al. (2023).
\newblock Anomalous citations detection in academic networks.
\newblock \emph{Artificial Intelligence Review}, 56.

\bibitem{heydari2025}
Ibrahim, H. et al. (2025).
\newblock Citation manipulation through citation mills and pre-print servers.
\newblock \emph{Scientific Reports}, 15.

\bibitem{ma2025graphad}
Qiao, H. et al. (2025).
\newblock Deep graph anomaly detection: A survey and new perspectives.
\newblock \emph{IEEE Transactions on Knowledge and Data Engineering}.

\bibitem{alkhalaf2025}
Abbonato, D. et al. (2025).
\newblock CheckIfExist: Detecting citation hallucinations in the era of AI-generated content.
\newblock \emph{arXiv preprint}, 2602.15871.

\bibitem{gravel2024}
Gravel, J., D'Amours-Gravel, M., \& Osmanlliu, E. (2023).
\newblock Learning to fake it: Limited responses and fabricated references provided by ChatGPT for medical questions.
\newblock \emph{Mayo Clinic Proceedings: Digital Health}, 1(3), 226--234.

\bibitem{wu2024llmdetect}
Wu, J. et al. (2024).
\newblock A survey on LLM-generated text detection: Necessity, methods, and future directions.
\newblock \emph{Computational Linguistics}, 51(1), 275--339.

\bibitem{wadden2022open}
Wadden, D. et al. (2022).
\newblock SciFact-Open: Towards open-domain scientific claim verification.
\newblock \emph{Findings of EMNLP}, 4719--4734.

\bibitem{chen2025copula}
Aich, A., Aich, A.B., \& Wade, B. (2025).
\newblock Deep copula classifier: Theory, consistency, and empirical evaluation.
\newblock \emph{arXiv preprint}, 2505.22997.

\bibitem{xu2025copula}
Aich, A. \& Hewage, S. (2025).
\newblock Copula based fusion of clinical and genomic machine learning risk scores.
\newblock \emph{arXiv preprint}, 2511.17605.

\bibitem{su2025ds}
Lu, Y. \& Zhu, W. (2025).
\newblock Soft likelihood-based multi-source fusion with reliability modeling under Dempster-Shafer theory.
\newblock \emph{Engineering Applications of Artificial Intelligence}, 148.

\bibitem{jiang2023ds}
Tang, Y. et al. (2023).
\newblock A new correlation belief function in Dempster-Shafer evidence theory and its application in classification.
\newblock \emph{Scientific Reports}, 13, 7609.

\bibitem{ding2025c3po}
Valkanas, A. et al. (2025).
\newblock C3PO: Optimized large language model cascades with probabilistic cost constraints for reasoning.
\newblock \emph{arXiv preprint}, 2511.07396.

\bibitem{straitouri2025}
Wang, T. \& Dobriban, E. (2025).
\newblock Optimal decision-making based on prediction sets.
\newblock \emph{arXiv preprint}, 2602.00989.

\bibitem{scancar2025ai}
Scancar, B. et al. (2025).
\newblock Revealing the paper mill iceberg: AI-based screening of cancer research publications.
\newblock \emph{bioRxiv preprint}, 2025.08.29.673016.

\bibitem{akritidis2025}
Fletcher, A.H.A. \& Stevenson, M. (2025).
\newblock Predicting retracted research: A dataset and machine learning approaches.
\newblock \emph{Research Integrity and Peer Review}, 10(1).

\bibitem{nakov2026}
Wang, H. et al. (2026).
\newblock Fact-checking with large language models via probabilistic certainty and consistency.
\newblock \emph{arXiv preprint}, 2601.02574.

\end{thebibliography}

\end{document}
